% =====================================================================
%  Insercao: a centralidade (baixa) de "tecnografia" na rede da tese
%
%  Observacao reflexiva sobre a analise de rede textual: o conceito que
%  nomeia a pratica tem centralidade estrutural baixa no grafo da propria
%  tese -- e por que isso e coerente com o argumento.
%
%  Duas formas de uso (escolha uma):
%   (a) nota de rodape -- prender a um "tecnografia" no corpo do texto,
%       de preferencia na seccao de analise de rede textual ou nas
%       Consideracoes finais.
%   (b) paragrafo -- inserir no corpo, na leitura da rede da tese inteira.
%
%  Numeros extraidos de infranodus/tese_network.json (rede da tese
%  inteira: Apresentacao + cap. 1--4 + Consideracoes finais).
%  Uso: % =====================================================================
%  Insercao: a centralidade (baixa) de "tecnografia" na rede da tese
%
%  Observacao reflexiva sobre a analise de rede textual: o conceito que
%  nomeia a pratica tem centralidade estrutural baixa no grafo da propria
%  tese -- e por que isso e coerente com o argumento.
%
%  Duas formas de uso (escolha uma):
%   (a) nota de rodape -- prender a um "tecnografia" no corpo do texto,
%       de preferencia na seccao de analise de rede textual ou nas
%       Consideracoes finais.
%   (b) paragrafo -- inserir no corpo, na leitura da rede da tese inteira.
%
%  Numeros extraidos de infranodus/tese_network.json (rede da tese
%  inteira: Apresentacao + cap. 1--4 + Consideracoes finais).
%  Uso: % =====================================================================
%  Insercao: a centralidade (baixa) de "tecnografia" na rede da tese
%
%  Observacao reflexiva sobre a analise de rede textual: o conceito que
%  nomeia a pratica tem centralidade estrutural baixa no grafo da propria
%  tese -- e por que isso e coerente com o argumento.
%
%  Duas formas de uso (escolha uma):
%   (a) nota de rodape -- prender a um "tecnografia" no corpo do texto,
%       de preferencia na seccao de analise de rede textual ou nas
%       Consideracoes finais.
%   (b) paragrafo -- inserir no corpo, na leitura da rede da tese inteira.
%
%  Numeros extraidos de infranodus/tese_network.json (rede da tese
%  inteira: Apresentacao + cap. 1--4 + Consideracoes finais).
%  Uso: % =====================================================================
%  Insercao: a centralidade (baixa) de "tecnografia" na rede da tese
%
%  Observacao reflexiva sobre a analise de rede textual: o conceito que
%  nomeia a pratica tem centralidade estrutural baixa no grafo da propria
%  tese -- e por que isso e coerente com o argumento.
%
%  Duas formas de uso (escolha uma):
%   (a) nota de rodape -- prender a um "tecnografia" no corpo do texto,
%       de preferencia na seccao de analise de rede textual ou nas
%       Consideracoes finais.
%   (b) paragrafo -- inserir no corpo, na leitura da rede da tese inteira.
%
%  Numeros extraidos de infranodus/tese_network.json (rede da tese
%  inteira: Apresentacao + cap. 1--4 + Consideracoes finais).
%  Uso: \input{docs/insercao-tecnografia-centralidade.tex}
% =====================================================================

% ---------------------------------------------------------------------
% (a) VERSAO NOTA DE RODAPE
% ---------------------------------------------------------------------
\footnote{%
Vale registrar um resultado que a rede me devolve sobre si mesma: o
conceito que nomeia esta pratica, \emph{tecnografia}, tem centralidade
estrutural baixa no grafo da propria tese. A palavra ocorre 62 vezes ---
contra 881 de \emph{rede} e 330 de \emph{etnografia} --- e sua
intermediacao (\emph{betweenness}) e nula: ela nao faz ponte entre
regioes do argumento, permanecendo como ponto terminal na comunidade do
metodo. Precisei fixar seu rotulo manualmente para que aparecesse no
mapa. Isso nao e falha do conceito, e a sua natureza: a rede mede
circulacao lexical e premia os descritores que espalho pelas descricoes
empiricas, ao passo que \emph{tecnografia} nao descreve, nomeia. Suas
associacoes mais fortes o confirmam --- \emph{conceito} (NPMI $0{,}40$),
\emph{nome} ($0{,}37$), \emph{palavra} ($0{,}29$), \emph{pratica},
\emph{descreve} ---, o vocabulario de batizar uma pratica, e nao o do
campo. O conceito enquadra a descricao sem se confundir com ela; que ele
emerja do processo analitico sem ganhar peso estatistico e, ele mesmo,
um resultado coerente com o que sustento.%
}

% ---------------------------------------------------------------------
% (b) VERSAO PARAGRAFO
% ---------------------------------------------------------------------
Um detalhe da rede me parece o mais eloquente de todos, e diz respeito ao
proprio nome deste trabalho. \emph{Tecnografia}, o conceito que nomeia a
pratica que descrevo, tem centralidade estrutural \textbf{baixa} no grafo
da tese inteira: ocorre 62 vezes, contra 881 de \emph{rede} --- o termo
mais central e a maior ponte do texto --- e 330 de \emph{etnografia}, e
sua intermediacao (\emph{betweenness}) e \textbf{nula}. Ela nao arbitra a
passagem entre as regioes do argumento; fica como ponto terminal na
comunidade do metodo, a tal ponto que precisei fixar seu rotulo a mao
para que nao desaparecesse entre os nos sem nome. Isso, porem, nao e uma
deficiencia do conceito, e a sua condicao. A rede mede \textbf{circulacao
lexical} --- recompensa as palavras que eu espalho ao longo das
descricoes ---, e \emph{tecnografia} nao e um descritor: e um gesto de
nomeacao. Suas co-ocorrencias mais fortes (NPMI) delatam essa natureza,
pois a palavra se associa a \emph{conceito}, \emph{nome}, \emph{palavra},
\emph{pratica} e \emph{descreve} --- o vocabulario de batizar uma
pratica, e nao o vocabulario do campo empirico. O conceito
\textbf{enquadra} a descricao sem se confundir com ela. Que ele tenha
emergido do processo analitico da tese sem ganhar peso estatistico
dentro dela e, assim, menos um paradoxo do que uma confirmacao: a
tecnografia e a lente, nao o objeto, e a rede --- fiel ao que a tese lhe
pede --- deixa isso a vista.

% =====================================================================

% ---------------------------------------------------------------------
% (a) VERSAO NOTA DE RODAPE
% ---------------------------------------------------------------------
\footnote{%
Vale registrar um resultado que a rede me devolve sobre si mesma: o
conceito que nomeia esta pratica, \emph{tecnografia}, tem centralidade
estrutural baixa no grafo da propria tese. A palavra ocorre 62 vezes ---
contra 881 de \emph{rede} e 330 de \emph{etnografia} --- e sua
intermediacao (\emph{betweenness}) e nula: ela nao faz ponte entre
regioes do argumento, permanecendo como ponto terminal na comunidade do
metodo. Precisei fixar seu rotulo manualmente para que aparecesse no
mapa. Isso nao e falha do conceito, e a sua natureza: a rede mede
circulacao lexical e premia os descritores que espalho pelas descricoes
empiricas, ao passo que \emph{tecnografia} nao descreve, nomeia. Suas
associacoes mais fortes o confirmam --- \emph{conceito} (NPMI $0{,}40$),
\emph{nome} ($0{,}37$), \emph{palavra} ($0{,}29$), \emph{pratica},
\emph{descreve} ---, o vocabulario de batizar uma pratica, e nao o do
campo. O conceito enquadra a descricao sem se confundir com ela; que ele
emerja do processo analitico sem ganhar peso estatistico e, ele mesmo,
um resultado coerente com o que sustento.%
}

% ---------------------------------------------------------------------
% (b) VERSAO PARAGRAFO
% ---------------------------------------------------------------------
Um detalhe da rede me parece o mais eloquente de todos, e diz respeito ao
proprio nome deste trabalho. \emph{Tecnografia}, o conceito que nomeia a
pratica que descrevo, tem centralidade estrutural \textbf{baixa} no grafo
da tese inteira: ocorre 62 vezes, contra 881 de \emph{rede} --- o termo
mais central e a maior ponte do texto --- e 330 de \emph{etnografia}, e
sua intermediacao (\emph{betweenness}) e \textbf{nula}. Ela nao arbitra a
passagem entre as regioes do argumento; fica como ponto terminal na
comunidade do metodo, a tal ponto que precisei fixar seu rotulo a mao
para que nao desaparecesse entre os nos sem nome. Isso, porem, nao e uma
deficiencia do conceito, e a sua condicao. A rede mede \textbf{circulacao
lexical} --- recompensa as palavras que eu espalho ao longo das
descricoes ---, e \emph{tecnografia} nao e um descritor: e um gesto de
nomeacao. Suas co-ocorrencias mais fortes (NPMI) delatam essa natureza,
pois a palavra se associa a \emph{conceito}, \emph{nome}, \emph{palavra},
\emph{pratica} e \emph{descreve} --- o vocabulario de batizar uma
pratica, e nao o vocabulario do campo empirico. O conceito
\textbf{enquadra} a descricao sem se confundir com ela. Que ele tenha
emergido do processo analitico da tese sem ganhar peso estatistico
dentro dela e, assim, menos um paradoxo do que uma confirmacao: a
tecnografia e a lente, nao o objeto, e a rede --- fiel ao que a tese lhe
pede --- deixa isso a vista.

% =====================================================================

% ---------------------------------------------------------------------
% (a) VERSAO NOTA DE RODAPE
% ---------------------------------------------------------------------
\footnote{%
Vale registrar um resultado que a rede me devolve sobre si mesma: o
conceito que nomeia esta pratica, \emph{tecnografia}, tem centralidade
estrutural baixa no grafo da propria tese. A palavra ocorre 62 vezes ---
contra 881 de \emph{rede} e 330 de \emph{etnografia} --- e sua
intermediacao (\emph{betweenness}) e nula: ela nao faz ponte entre
regioes do argumento, permanecendo como ponto terminal na comunidade do
metodo. Precisei fixar seu rotulo manualmente para que aparecesse no
mapa. Isso nao e falha do conceito, e a sua natureza: a rede mede
circulacao lexical e premia os descritores que espalho pelas descricoes
empiricas, ao passo que \emph{tecnografia} nao descreve, nomeia. Suas
associacoes mais fortes o confirmam --- \emph{conceito} (NPMI $0{,}40$),
\emph{nome} ($0{,}37$), \emph{palavra} ($0{,}29$), \emph{pratica},
\emph{descreve} ---, o vocabulario de batizar uma pratica, e nao o do
campo. O conceito enquadra a descricao sem se confundir com ela; que ele
emerja do processo analitico sem ganhar peso estatistico e, ele mesmo,
um resultado coerente com o que sustento.%
}

% ---------------------------------------------------------------------
% (b) VERSAO PARAGRAFO
% ---------------------------------------------------------------------
Um detalhe da rede me parece o mais eloquente de todos, e diz respeito ao
proprio nome deste trabalho. \emph{Tecnografia}, o conceito que nomeia a
pratica que descrevo, tem centralidade estrutural \textbf{baixa} no grafo
da tese inteira: ocorre 62 vezes, contra 881 de \emph{rede} --- o termo
mais central e a maior ponte do texto --- e 330 de \emph{etnografia}, e
sua intermediacao (\emph{betweenness}) e \textbf{nula}. Ela nao arbitra a
passagem entre as regioes do argumento; fica como ponto terminal na
comunidade do metodo, a tal ponto que precisei fixar seu rotulo a mao
para que nao desaparecesse entre os nos sem nome. Isso, porem, nao e uma
deficiencia do conceito, e a sua condicao. A rede mede \textbf{circulacao
lexical} --- recompensa as palavras que eu espalho ao longo das
descricoes ---, e \emph{tecnografia} nao e um descritor: e um gesto de
nomeacao. Suas co-ocorrencias mais fortes (NPMI) delatam essa natureza,
pois a palavra se associa a \emph{conceito}, \emph{nome}, \emph{palavra},
\emph{pratica} e \emph{descreve} --- o vocabulario de batizar uma
pratica, e nao o vocabulario do campo empirico. O conceito
\textbf{enquadra} a descricao sem se confundir com ela. Que ele tenha
emergido do processo analitico da tese sem ganhar peso estatistico
dentro dela e, assim, menos um paradoxo do que uma confirmacao: a
tecnografia e a lente, nao o objeto, e a rede --- fiel ao que a tese lhe
pede --- deixa isso a vista.

% =====================================================================

% ---------------------------------------------------------------------
% (a) VERSAO NOTA DE RODAPE
% ---------------------------------------------------------------------
\footnote{%
Vale registrar um resultado que a rede me devolve sobre si mesma: o
conceito que nomeia esta pratica, \emph{tecnografia}, tem centralidade
estrutural baixa no grafo da propria tese. A palavra ocorre 62 vezes ---
contra 881 de \emph{rede} e 330 de \emph{etnografia} --- e sua
intermediacao (\emph{betweenness}) e nula: ela nao faz ponte entre
regioes do argumento, permanecendo como ponto terminal na comunidade do
metodo. Precisei fixar seu rotulo manualmente para que aparecesse no
mapa. Isso nao e falha do conceito, e a sua natureza: a rede mede
circulacao lexical e premia os descritores que espalho pelas descricoes
empiricas, ao passo que \emph{tecnografia} nao descreve, nomeia. Suas
associacoes mais fortes o confirmam --- \emph{conceito} (NPMI $0{,}40$),
\emph{nome} ($0{,}37$), \emph{palavra} ($0{,}29$), \emph{pratica},
\emph{descreve} ---, o vocabulario de batizar uma pratica, e nao o do
campo. O conceito enquadra a descricao sem se confundir com ela; que ele
emerja do processo analitico sem ganhar peso estatistico e, ele mesmo,
um resultado coerente com o que sustento.%
}

% ---------------------------------------------------------------------
% (b) VERSAO PARAGRAFO
% ---------------------------------------------------------------------
Um detalhe da rede me parece o mais eloquente de todos, e diz respeito ao
proprio nome deste trabalho. \emph{Tecnografia}, o conceito que nomeia a
pratica que descrevo, tem centralidade estrutural \textbf{baixa} no grafo
da tese inteira: ocorre 62 vezes, contra 881 de \emph{rede} --- o termo
mais central e a maior ponte do texto --- e 330 de \emph{etnografia}, e
sua intermediacao (\emph{betweenness}) e \textbf{nula}. Ela nao arbitra a
passagem entre as regioes do argumento; fica como ponto terminal na
comunidade do metodo, a tal ponto que precisei fixar seu rotulo a mao
para que nao desaparecesse entre os nos sem nome. Isso, porem, nao e uma
deficiencia do conceito, e a sua condicao. A rede mede \textbf{circulacao
lexical} --- recompensa as palavras que eu espalho ao longo das
descricoes ---, e \emph{tecnografia} nao e um descritor: e um gesto de
nomeacao. Suas co-ocorrencias mais fortes (NPMI) delatam essa natureza,
pois a palavra se associa a \emph{conceito}, \emph{nome}, \emph{palavra},
\emph{pratica} e \emph{descreve} --- o vocabulario de batizar uma
pratica, e nao o vocabulario do campo empirico. O conceito
\textbf{enquadra} a descricao sem se confundir com ela. Que ele tenha
emergido do processo analitico da tese sem ganhar peso estatistico
dentro dela e, assim, menos um paradoxo do que uma confirmacao: a
tecnografia e a lente, nao o objeto, e a rede --- fiel ao que a tese lhe
pede --- deixa isso a vista.
